%% get rid of autoindent   :setl noai nocin nosi inde=
%%
%% 1. pdflatex main
%% 2. bibtex main
%% 3. pdflatex main
%% 3. pdflatex main
%\documentclass[review]{elsarticle}
\documentclass[]{article}
%\documentclass[preprint]{elsarticle}

\usepackage[fleqn]{amsmath} 
\usepackage{amssymb}
\usepackage{braket}
\usepackage{amsmath}
\usepackage{graphicx}
\usepackage{tabularx} 
\usepackage{footnote}
\usepackage{slashed}
%%\usepackage{titlesec}
\usepackage{comment}

%\journal{Journal of Computational Physics}

\newcommand{\mb}{\mathbf}
\newcommand{\tn}{\textnormal}
\newcommand{\unit}[1]{\ensuremath{\, \mathrm{#1}}}

\newcommand{\bmz}{\mbox{\boldmath $z$\unboldmath}}
\newcommand{\bmr}{\mbox{\boldmath $r$\unboldmath}}
\newcommand{\bmb}{\mbox{\boldmath $b$\unboldmath}}
\newcommand{\bms}{\mbox{\boldmath $s$\unboldmath}}
\newcommand{\bmg}{\mbox{\boldmath $g$\unboldmath}}
\newcommand{\bmq}{\mbox{\boldmath $q$\unboldmath}}
\newcommand{\bmU}{\mbox{\boldmath $U$\unboldmath}}
\newcommand{\bmu}{\mbox{\boldmath $u$\unboldmath}}
\newcommand{\bmv}{\mbox{\boldmath $v$\unboldmath}}
\newcommand{\bmw}{\mbox{\boldmath $w$\unboldmath}}
\newcommand{\bmm}{\mbox{\boldmath $m$\unboldmath}}
\newcommand{\bmi}{\mbox{\boldmath $i$\unboldmath}}
\newcommand{\bmj}{\mbox{\boldmath $j$\unboldmath}}
\newcommand{\bmk}{\mbox{\boldmath $k$\unboldmath}}
\newcommand{\bmx}{\mbox{\boldmath $x$\unboldmath}}
\newcommand{\bmX}{\mbox{\boldmath $X$\unboldmath}}
\newcommand{\bmH}{\mbox{\boldmath $H$\unboldmath}}
\newcommand{\bmy}{\mbox{\boldmath $y$\unboldmath}}
\newcommand{\bmn}{\mbox{\boldmath $n$\unboldmath}}
\newcommand{\bmf}{\mbox{\boldmath $f$\unboldmath}}
\newcommand{\bmF}{\mbox{\boldmath $F$\unboldmath}}
\newcommand{\bmS}{\mbox{\boldmath $S$\unboldmath}}
\newcommand{\bmG}{\mbox{\boldmath $G$\unboldmath}}
\newcommand{\bmV}{\mbox{\boldmath $V$\unboldmath}}
\newcommand{\bmI}{\mbox{\boldmath $I$\unboldmath}}
\newcommand{\bmD}{\mbox{\boldmath $D$\unboldmath}}
\newcommand{\bmnu}{\mbox{\boldmath $\nu$\unboldmath}}
\newcommand{\bmalpha}{\mbox{\boldmath $\alpha$\unboldmath}}

\newcommand{\DefTen}{\mathbb{D}}
\newcommand{\EyeTen}{\mathbb{I}}

\DeclareMathOperator{\Ei}{Ei}

%%\bibliographystyle{elsarticle-num}
\bibliographystyle{acm}
%%%%%%%%%%%%%%%%%%%%%%%

\title{Review article: Numerical Methods for Engineering Quantum Computers}

%%!!format authors properly
%\author{
%	LastName1, FirstName1\\
%	\texttt{first1.last1@xxxxx.com}
%	\and
%	LastName2, FirstName2\\
%	\texttt{first2.last2@xxxxx.com}
%}
\author{
  Liu, Shu \\
  Mathematics, Florida State University
  \and
  Georgiadou, Antigoni \\
  Oak Ridge National Laboratory
  \and
  Dong, Spancer \\
  Oak Ridge National Laboratory
  \and
  Kubischta, Eric \\
  Mathematics, Florida State University
  \and
  Melton, Daniel \\
  Mathematics, Florida State University
  \and
  Guenther, Stefanie \\
  Lawrence Livermore National Laboratory
  \and
  Jin, Shi \\
  Mathematics, Shanghai Jiao Tong University
  \and
  Lin, Lin \\
  Mathematics, U.C. Berkeley
  \and
  Lotshaw, Philip \\
  Oak Ridge National Laboratory
  \and
  Petersson, Anders \\
  Lawrence Livermore National Laboratory
  \and
  Sussman, Mark \\
  Mathematics, Florida State University
}

\begin{document}
\maketitle

ABSTRACT:
In the article by Herrmann et al, ``Quantum utility – definition and 
assessment of a practical quantum advantage,'' the ``Application
Readiness Level'' (ARL) is defined: (ARL1) ``concept with unknown potential,'' 
(ARL2) ``beneficial in small idealized systems,'' (ARL3) ``utility indicated by 
theory or resource estimations,'' (ARL4) ``simulated utility demonstration,''
and (ARL5) ``utility demonstration on quantum hardware.''  At the present, 
practitioners have reached ARL3 with the Variational Quantum Eigenvalue
(VQE) algorithm.  This seems to indicate that there is still much work
to be done in order to reach ARL4 or ARL5.
In this review article, the latest activities from the
computational ``digital twin'' perspective are reported on in which scientists
are working hard to bring practitioners out of the Noisy Intermediate
Scale Quantum (NISQ) era into an era of ``Quantum supremacy.'' In
otherwords this review article discusses 
digital twin activities that help to advance from ARL3 to 
ARL4, or even ARL5.  The three main categories in this review article are
(1) Quantum algorithms for benchmarking quantum computers (to include
emulator algorithms like Qiskit), (2) Optimization, Control, and inverse
problem algorithms for Schrodinger's equation, and (3) Error correction
algorithms and algorithms for optimizing the ``surface topology'' of
logical qubits.

\par\noindent
The motivation for this review article is the fact that Moore's law is no 
longer applicable\cite{moore2021cramming}.
As quoted from Shalf\cite{shalf2020future},
``Moore’s Law [1] is a techno-economic model that has enabled 
the IT industry to double the performance and 
functionality of digital electronics roughly 
every 2 years within a fixed cost, power and area. 
This expectation has led to a relatively stable ecosystem (e.g. 
electronic design automation tools, compilers, simulators and 
emulators) built around general-purpose processor technologies, 
such as the ×86, ARM and Power instruction set architectures. 
However, within a decade, the technological underpinnings 
for the process that Gordon Moore described will come to an end, 
as lithography gets down to atomic scale. At that point, 
it will be feasible to create lithographically produced devices 
with dimensions nearing atomic scale, where a dozen or fewer 
silicon atoms are present across critical device features, and 
will therefore represent a practical limit for implementing logic 
gates for digital computing [2]. Indeed, the ITRS (International 
Technology Roadmap for Semiconductors), which has tracked the 
historical improvements over the past 30 years, has projected no 
improvements beyond 2021, as shown in figure 1, and 
subsequently disbanded, having no further purpose. 
The classical technological driver that has underpinned Moore’s Law 
for the past 50 years is failing [3] and is anticipated to flatten by 
2025, as shown in figure 2. Evolving technology in the absence of 
Moore’s Law will require an investment now in computer 
architecture and the basic sciences (including materials science), 
to study candidate replacement materials and alternative 
device physics to foster continued technology scaling.'' 
\par\noindent
References one through three in the above quote 
correspond to the following references respectively: 
\cite{moore2021cramming}\cite{mack2015multiple}\cite{markov2014limits}.
\par\noindent
As further motivation for this review article,
in Figure 3 of the article by Shalf\cite{shalf2020future}, a 
roadmap is provided for possible paths forward when the 
density of circuits on classical computers exceeds a critical 
value.  There are three categories: (a) roadmap for the next ten years, 
(b) 20 years, and (c) ``Decades beyond exascale,'' 
``New Models of Computation.'' For category (c), Shalf\cite{shalf2020future}
refers to the following
prospective technology:
(i) approximate computing, (ii) adiabatic reversible, 
(iii) Analog, (iv) Neuromorhic, and (v) quantum.
\par\noindent
Of the ``new models of computation'' listed in 
Shalf's article\cite{shalf2020future}, this review article will address
the ``Numerical Methods for Engineering Quantum Computers'' 
aspects associated with the emerging ``quantum computing'' paradigm.  
As outlined by
Gamble\cite{gamble2019quantum}, the development of reliable quantum
computers will have a tranformative effect on computer technology.  
That being said, right now we are in the 
``Noisy Intermediate-Scale Quantum''\cite{callison2022hybrid} (NISQ) era.  
The purpose of this review article is to bring to the 
forefront the current research being done, from the 
computational perspective, in order to move beyond the NISQ era.

\section{Quantum algorithms and Quantum emulators
  for benchmarking quantum computers}

\par\noindent
Quantum Emulators:
\begin{enumerate}
\item (IBM) ``Quantum Computing with Qiskit'' (2024)\cite{javadi2024quantum}
\item ``Intel Quantum Simulator: a cloud-ready high-performance
 simulator of quantum circuits'' (2020)\cite{Guerreschi_Intel_Quantum_Simulator_2020}
\item (Microsoft) ``LIQUid: A Software design architecture 
 and domain specific language
 for quantum computing'' (2014)\cite{1402.4467}
\item ``QuESTlink—Mathematica embiggened by a hardware-optimised quantum emulator'' (2020)\cite{Jones_2020}
\item (Matlab) ``QCLAB: A Matlab Toolbox for Quantum Computing'' 
 (2025)\cite{keip2025qclabmatlabtoolboxquantum}
\item ``PennyLane-Lightning MPI: A massively scalable quantum circuit simulator based on distributed computing in CPU clusters'' (2025)\cite{kang2025pennylane}
\item (Google) ``qsim''(2020)\cite{quantum_ai_team_and_collaborators_2020_4023103}
\item ``Numerical recipes in quantum information theory and quantum computing: an adventure in FORTRAN 90''(2021)\cite{ramkarthik2021numerical}
\item Nguyen et al(2022)\cite{nguyen2022tensor}
``Tensor network quantum virtual machine for simulating quantum circuits at exascale''
\item Zhang et al\cite{zhang2024quantum}, ``Quantum circuit simulation with fast tensor decision diagram''
\end{enumerate}

\par\noindent
Quantum Algorithms for benchmarking Quantum Computers:
\begin{enumerate}
\item Callison and Chancellor (2022)\cite{callison2022hybrid} ``Hybrid quantum-classical algorithms in the noisy intermediate-scale quantum era and beyond''
\item S. Jin, N. Liu, and Y. Yu (2024)\cite{jin2024quantum}, ``Quantum Circuits for the heat equation with physical boundary conditions via Schrodingerisation''
\item 
Cerezo et al (2021)\cite{cerezo2021variational} ``Variational quantum algorithms''
\item
Lotshaw et al (2021)\cite{doi:10.1098/rsta.2021.0414} ``Simulations of frustrated Ising Hamiltonians using quantum approximate optimization''
\item Nielsen and Chuang (2010)\cite{nielsen2010quantum}
``Quantum computation and quantum information'' (quantum Fourier Transform)
\item Ding and Lin (2023)\cite{ding2023simultaneous}
``Simultaneous estimation of multiple eigenvalues with short-depth quantum circuit on early fault-tolerant quantum computers''
\item Biamonte et al (2017)\cite{biamonte2017quantum}
``Quantum machine learning''
\item Benedetti et al (2019)\cite{Benedetti_2019}
``Parameterized quantum circuits as machine learning models''
\item Khait et al(2023)\cite{khait2023variational}
''Variational quantum eigensolvers in the era of distributed quantum computers''
\item Sawaya et al (2023)\cite{10313872},
  ``HamLib: A Library of Hamiltonians for Benchmarking Quantum Algorithms and Hardware''
\item Georgadou et al(2024)\cite{gopalakrishnan2024solving}
	Hele-Shaw simulation on quantum computer.
\item Kim et al (2023)\cite{kim2023evidence}, ``Evidence for the utility of quantum computing before fault tolerance''
\item Zoufal et al (2019)\cite{zoufal2019quantum}, ``Quantum generative adversarial networks for learning and loading random distributions''
\item Nghiem and Tzu-Chieh (2023)\cite{nghiem2023quantum}, ``Quantum algorithm for estimating largest eigenvalues''
\item Cadi et al (2024)\cite{cadi2024folded}, ``Folded spectrum vqe: A quantum computing method for the calculation of molecular excited states''
\item Tilly et al (2022)\cite{tilly2022variational}, ``The variational quantum eigensolver: a review of methods and best practices''
\item Qing and Xie (2023)\cite{qing2023usevqecalculateground}, ``Use VQE to calculate the ground energy of hydrogen molecules on IBM Quantum''
\item Ryu et al (2021)\cite{VQSVD}, ``A Variational Quantum Algorithm for Ordered SVD''
\item Shirakawa et al (2024)\cite{PhysRevResearch.6.043008}, ``Automatic quantum circuit encoding of a given arbitrary quantum state''
\item Harrow et al (2009)\cite{harrow2009quantum}, ``Quantum algorithm for linear systems of equations''
\item Daribayev et al (2023)\cite{daribayev2023implementation}, 
  ``Implementation of the hhl algorithm for solving the poisson equation on quantum simulators.''
\item Bravo-Prieto et al (2023)\cite{bravo2023variational},
  ''Variational quantum linear solver''
\end{enumerate}

\begin{comment}
	1. PDE constrained optimization and control. QUANDARY
	  (PDE constraint is Schrodinger equation)
	  Anders Petersson and Stefanie Guenther
	2. Quantum algorithms for benchmarking quantum computers
	   Shi Jin, Antigoni Georgiadou (discuss emulators Qiskit,
           Penny Lane, ...; when discussing emulators, also discuss how they
	   emulate the noise of a real quantum computer) 
	   VQE, Shor's prime factorization algorithm,
           Quantum Fourier Transform
	3. Quantum error correction algorithms
	   Eric Kubischta
\end{comment}

\section{Quantum Algorithms for benchmarking Quantum Computers}

\verb=https://www.ibm.com/quantum/qiskit=
\begin{comment}
mkdir qiskit
cd qiskit
pip3 install qiskit
python3 -m venv venv_sussman
source venv_sussman/bin/activate (this should be done every session)
python3 -m pip install qiskit
pip3 install qiskit-aer
pip3 install matplotlib
pip3 install qiskit-ibm-runtime
python3 hadamard_gate.py
\end{comment}

Discuss the latest technologies for Quantum Computer Emulators.

Examples: (1)VQE using slater determinant form, only valence electrons, 
(2) SVD (SVT, self consitent field theory), 
(3) Shor's prime number factorization algorithm, (4) heat 
equation, and (5) Quantum Fourier Transform.

\begin{itemize}
\item Jin et al\cite{doi:10.1137/23M1563451}
\item Beck et al\cite{beck2024integrating}
\end{itemize}

\section{Quantum Algorithms with error correction logic}

\begin{itemize}
\item Shor's algorithm for bit flip and phase flip errors.  e.g. given
 bit flip probability $p$, what is the probability of a bit flip 
 error occurring with and without error correction?
\item Kubischta, Eric and Teixeira, Ian\cite{kubischta2023family}
\end{itemize}

\begin{itemize}
\item Preskill's notes chapter 7.
\item Eastin and Knill: A transversal quantum gate does not spread
  errors between physical qubits. ``if there are $n_{e}$ errors in the
  physical qubits representing the logical system, then after the 
  application of a transversal gate, there will still be $n_{e}$ errors.''
  i.e., if the number of errors is below the ``correctable'' threshold $t$
  prior to the application of a transversal gate, it will still be 
  correctable after the application of a tranversal gate.
\item Kubischta and Teixeira (a quantum code for making the 2I gates
  transversal (i.e. product gates))
\item Ide et al\cite{ide2025fault} (surface codes)
\end{itemize}

\section{Control and optimization of Quantum Computers using 
  numerical model surrogates in place of the actual quantum computer}

\begin{itemize}
\item Günther et al\cite{9651392}
\item Fei et al\cite{fei2025switching}
\end{itemize}

\subsection{Control of open quantum systems}

The Lindblad master equation for Open quantum systems\cite{lidar2019lecture} is
\begin{eqnarray}
\dot{\rho}(t)=-i[H(t),\rho(t)]+\ell(\rho(t)) \hspace{10pt} 0\le t\le T
\end{eqnarray}
\begin{eqnarray}
\ell(\rho)=\sum_{q=1}^{Q}\sum_{l=1}^{2} \ell_{lq}\rho\ell^{\dag}_{lq}-
  \frac{1}{2}(\ell^{\dag}_{lq}\ell_{lq}\rho+\rho\ell^{\dag}_{lq}\ell_{lq})
\end{eqnarray}
\begin{eqnarray}
N=\Pi_{q=1}^{Q}n_{q}
\end{eqnarray}
\begin{eqnarray}
\rho(t)\in\mathbb{C}^{N\times N}
\end{eqnarray}
From \cite{lidar2019lecture}, given the pure state ensemble,
\begin{eqnarray}
\{q_{i},\ket\psi_{i}\}_{i=1}^{N},
\end{eqnarray}
the density operator is
\begin{eqnarray}
\rho=\sum_{i}q_{i}\ket\psi_{i}\bra\psi_{i}.
\end{eqnarray}
\begin{eqnarray}
H=H_{sys}+H_{c}(t)
\end{eqnarray}
The control Hamiltonian has the form
\begin{eqnarray}
H_{c}(t)=\sum_{q=1}^{Q}d_{q}(\alpha_{q},t)a_{q}+
  d_{q}(\alpha_{q},t)^{\ast}(t)a_{q}^{\dag}.
\end{eqnarray}
The objective function has the form
\begin{eqnarray}
{\cal G}(\bmalpha)=\frac{1}{M}\sum_{i=1}^{M} \left( 
\beta_{i} J(\rho_{i}^{target},\rho_{i}(T))+
\gamma_{1}\int_{0}^{T}\beta(t)J(\rho_{i}^{target},\rho_{i}(t))dt \right)+
\gamma_{2} ||\bmalpha||_{2}^{2}/2
\end{eqnarray}
In Quandary the following choices for $J$ are\cite{9651392}
\begin{eqnarray}
J_{Frob}(\rho^{target},\rho)=||\rho^{target}-\rho||^{2}_{F}/2
\end{eqnarray}
\begin{eqnarray}
J_{Trace}(\rho^{target},\rho)=1-Tr((\rho^{target})^{\dag}\rho)
\end{eqnarray}
\begin{eqnarray}
J_{N_{m}}(\rho)=Tr(N_{m}\rho)
\end{eqnarray}

\subsection{Switching time optimization Control of binary quantum systems}

\begin{eqnarray}
\dot{X}(t)=-iH(t)X(t) \hspace{10pt} 0\le t\le T
\end{eqnarray}
\begin{eqnarray}
H(t)=H^{(0)}+\sum_{j=1}^{J} u_{j}(t)H^{(j)}  \hspace{10pt} 0\le t\le T
\end{eqnarray}
\begin{eqnarray}
\sum_{j=1}^{J}u_{j}(t)=1 \hspace{10pt} 0\le t\le T
\end{eqnarray}
\begin{eqnarray}
u_{j}(t) \in \left\{ 0,1 \right\} \hspace{10pt} j=1,\ldots,J
\end{eqnarray}
Cost function example ($Q=2$):
\begin{eqnarray}
{\cal G}(X[u_{1}(t),\ldots,u_{J}(t)](T))=1-\frac{1}{4}|\mbox{Trace}\left\{X_{CNOT}^{\dag}X(T)\right\}|
\end{eqnarray}

\begin{comment}

Explains how the anticipated contribution of the special issue will advance understanding in this area;

\end{comment}
\begin{comment}

At the present, the research activity associated with ``Numerical Methods for Engineering Quantum Computers''
is spread out over many journals: IEEE Journals, 
``Quantum Information Processing,'' 
``Quantum Science and Technology,'' 
``Quantum,'' 
``ACM Transactions on Quantum Computing,''
``SIAM review,''
``Journal of Computational Physics,''
``Journal of Scientific Computing,''
``Physical Review A,'' 
``Journal of Chemical Physics,'' 
``Journal of Physics: Condensed Matter,''
``Physical Review A,''
``Physical Review Letters,''
``Physical Review Research,''
``Nature,''
``Nature Communications,''
``Nature Physics,''
``Nature Photonics,''
``Nature Chemistry,''
``PRX Quantum,''
``AVS Quantum Science,''
``Philosophical Transactions of the Royal Society A: Mathematical, Physical and Engineering Sciences.''

It is the intention of this special issue to have information pertaining
to ``Numerical Methods for Engineering Quantum Computers'' in one accessble issue.
\end{comment}

\begin{comment}

Identifies papers and authors for possible inclusion in the special issue, with a brief description of each paper. (These papers do not need to have been written at this time, although it might be the case that work is already in progress.);

\end{comment}

%Oak Ridge: Travis Humble, Stephen Jesse, Benjamin Lawrie, 
%Gilles Buchs, Rob Moore, Zac Ward
%FSU: Wei Guo, Lukasz Dusanowski
\begin{comment}

The following is a list of authors and the articles they
have previously published which motivate inviting them to 
contribute to the proposed special issue on ``Numerical Methods for Engineering Quantum Computers.''

\begin{itemize}
\item Callison and Chancellor\cite{callison2022hybrid} 
 ``Hybrid quantum-classical algorithms in the noisy intermediate-scale quantum era and beyond''
\item Muqeet et al\cite{muqeet2024machine}
``A Machine Learning-Based Error Mitigation Approach for Reliable Software Development on IBM’s Quantum Computers''
\item Anthony-Petersen et al\cite{anthony2024stress}
``A stress-induced source of phonon bursts and quasiparticle poisoning''
\item Tuysuz et al\cite{tuysuz2024learning}
  ``Learning to generate high-dimensional distributions with low-dimensional quantum Boltzmann machines''
\item W. Bao, Z. Chang, X. Zhao\cite{BAO2025113486}, ``Computing ground states of Bose-Einstein condensation by normalized deep neural network''
\item W. Bao, S. Jin, and P. Markowich\cite{BAO2002487}
``On Time-Splitting Spectral Approximations for the Schrödinger Equation in the Semiclassical Regime''
\item S. Jin, X. Li, N. Liu, and Y. Yu\cite{doi:10.1137/23M1563451},
 ``Quantum Simulation for Quantum Dynamics with Artificial Boundary 
Conditions''
\item S. Jin, N. Liu, and Y. Yu\cite{jin2024quantum}, ``Quantum Circuits for the heat equation with physical boundary conditions via Schrodingerisation''
\item S. Jin, H. Liu, S. Osher, R. Tsai\cite{jin2005computing}, ``Computing multivalued physical observables for the semiclassical limit of the Schr{\"o}dinger equation''
\item
Jiequn Han and Linfeng Zhang and Weinan E\cite{HAN2019108929}
``Solving many-electron Schrödinger equation using deep neural networks''
\item
Bernien et al\cite{bernien2017probing}
``Probing many-body dynamics on a 51-atom quantum simulator''
\item
Zhang et al\cite{zhang2017observation}
``Observation of a many-body dynamical phase transition with a 53-qubit quantum simulator''
\item 
Cerezo et al\cite{cerezo2021variational}
``Variational quantum algorithms''
\item
Anshu et al\cite{anshu2021sample}
``Sample-efficient learning of interacting quantum systems''
\item
Toshiaki Kanai, Dafei Jin, and Wei Guo\cite{PhysRevLett.132.250603}
``Single-Electron Qubits Based on Quantum Ring States on Solid Neon Surface''
\item
Hermann et al\cite{hermann2020deep}
``Deep-neural-network solution of the electronic Schr{\"o}dinger equation''
\item
Ilin and Arad\cite{ilin2024dissipativevariationalquantumalgorithms}
``Dissipative variational quantum algorithms for Gibbs state preparation''
\item
Somma et al\cite{somma2024shadowhamiltoniansimulation}
``Shadow Hamiltonian Simulation''
\item 
Lotshaw et al\cite{PhysRevA.110.L030803}
``Exactly solvable model of light-scattering errors in quantum simulations with metastable trapped-ion qubits''
\item
Lotshaw et al\cite{PhysRevA.107.062406}
``Modeling noise in global M\o{}lmer-S\o{}rensen interactions applied to quantum approximate optimization''
\item
Lotshaw et al\cite{doi:10.1098/rsta.2021.0414}
``Simulations of frustrated Ising Hamiltonians using quantum approximate optimization''
\item Friesen et al\cite{PhysRevB.67.121301}
  ``Practical design and simulation of silicon-based quantum-dot qubits''
\item Kai Jiang et al\cite{jiang2024high}
``High-accuracy numerical methods and convergence analysis for Schr{\"o}dinger equation with incommensurate potentials''
\item Lin and Lu\cite{lin2019mathematical}
``A mathematical introduction to electronic structure theory''
\item Nielsen and Chuang\cite{nielsen2010quantum}
``Quantum computation and quantum information''
\item Ding and Lin\cite{ding2023simultaneous}
``Simultaneous estimation of multiple eigenvalues with short-depth quantum circuit on early fault-tolerant quantum computers''
\item Ding et al\cite{ding2024random}
  ``Random coordinate descent: A simple alternative for optimizing parameterized quantum circuits''
\item Liu and Lin\cite{LIU2024113213}
``Dense outputs from quantum simulations''
\item Lin, Saad, and Yang\cite{lin2016approximating}
``Approximating spectral densities of large matrices''
\item Kubischta and Teixeira\cite{kubischta2023family}
``Family of quantum codes with exotic transversal gates''
\item Chen et al\cite{10628380}
``Quantum-Classical-Quantum Workflow in Quantum-HPC Middleware with GPU Acceleration''
\item Biamonte et al\cite{biamonte2017quantum}
``Quantum machine learning''
\item Benedetti et al\cite{Benedetti_2019}
``Parameterized quantum circuits as machine learning models''
\item Jin-Guo Liu et al\cite{PhysRevResearch.1.023025}
''Variational quantum eigensolver with fewer qubits''
\item Saurabh et al\cite{saurabh2023conceptual}
``A conceptual architecture for a quantum-hpc middleware''
\item Uvarov et al\cite{uvarov2020machine}
``Machine learning phase transitions with a quantum processor''
\item Khait et al\cite{khait2023variational}
''Variational quantum eigensolvers in the era of distributed quantum computers''
\item Huang et al\cite{huang2020predicting}
``Predicting many properties of a quantum system from very few measurements''
\item Cong et al\cite{cong2019quantum}
``Quantum convolutional neural networks''
\item Parekh et al\cite{parekh2021quantum}
``Quantum algorithms and simulation for parallel and distributed quantum computing''
\item Vallero et al\cite{vallero2024state}
``State of practice: evaluating GPU performance of state vector and tensor network methods''
\item Nguyen et al\cite{nguyen2022tensor}
``Tensor network quantum virtual machine for simulating quantum circuits at exascale''
\item Liu et al\cite{liu2024training}
``Training classical neural networks by quantum machine learning''
\item Bayraktar et al\cite{10313722}
``cuQuantum SDK: A High-Performance Library for Accelerating Quantum Science''
\item Andrade et al\cite{Andrade_2022}
``Engineering an effective three-spin Hamiltonian in trapped-ion systems for applications in quantum simulation''
%%\item https://github.com/lin-lin/Quantum290
\item Wenhao He et al
 \cite{he2024efficientoptimalcontrolopen},
 ``Efficient Optimal Control of Open Quantum Systems''
\item Alexander Nusseler et al
\cite{PhysRevB.101.155134},
``Efficient simulation of open quantum systems coupled to a fermionic bath''
\item Selsto and Kvaal
\cite{Selsto2010},
``Absorbing boundary conditions for dynamical many-body quantum systems''
\item Sawaya et al\cite{10313872},
  ``HamLib: A Library of Hamiltonians for Benchmarking Quantum Algorithms and Hardware''
\item Symeon Grivopoulos\cite{grivopoulos2005optimal},
 ``Optimal control of quantum systems''
\item Theisen and Stamm\cite{doi:10.1137/23M161848X},
  ''A Scalable Two-Level Domain Decomposition Eigensolver for Periodic Schrödinger Eigenstates in Anisotropically Expanding Domains''
\item Guenther, Petersson, and DuBois\cite{9651392},
  ``Quandary: An open-source C++ package for high-performance optimal control of open quantum systems''
\item Guenther and Petersson\cite{gunther2023practical},
``A practical approach to determine minimal quantum gate durations using amplitude-bounded quantum controls''
\item Lidar\cite{lidar2019lecture},
``Lecture notes on the theory of open quantum systems''
\item Hangleiter et al\cite{hangleiter2024robustly},
``Robustly learning the Hamiltonian dynamics of a superconducting quantum processor''
\item Beer et al\cite{beer2020training},
``Training deep quantum neural networks''
\item Lalanne et al\cite{LalanneETAL},
``Light Interaction with Photonic and Plasmonic Resonances''
\item Shukrinov et al\cite{shukrinov2016modeling},
``Modeling of LC-shunted intrinsic Josephson junctions in high-Tc superconductors''
\item Georgadou et al\cite{gopalakrishnan2024solving}
	Hele-Shaw simulation on quantum computer.
\item Rupert Klein\cite{riedel2023wavetrain} WaveTrain: A Python package for numerical quantum mechanics of chain-like systems based on tensor trains
\item {\verb=psiquantum.com=} \cite{bartolucci2023fusion}  ``Fusion based quantum computation,'' Bartolucci et al
\item Roy Goodman et al\cite{goodman2025qglab}, ``QGLAB: A MATLAB PACKAGE FOR COMPUTATIONS ON QUANTUM GRAPHS''
\item Campolattaro, Alfonso A ``Generalized Maxwell equations and quantum mechanics. I. Dirac equation for the free electron''
\item Brannick et al, ``Least-Squares Finite Element Methods for Quantum Electrodynamics''
\item Roth and Chew, ``Macroscopic Circuit Quantum Electrodynamics: A New Look Toward Developing Full-Wave Numerical Models''
\item Devoret et al, ``Quantum fluctuations in electrical circuits''
\end{itemize}
\end{comment}

\begin{comment}

Indicates the timeframe in which the special issue could be produced (to include paper writing, reviewing and submission of final copy to the journal) assuming the proposal is accepted;
Includes a short biography of all authors and guest editors;
Indicates any special timing, associated events, funding support, partnerships or other links or relationships which could influence the development of the issue;
Provides any further information which you feel is relevant.
A special issue normally contains between five and 20 full-length articles, in addition to an editorial written by the special issue organizers. Because it is highly unlikely that all articles submitted for potential inclusion in a special issue will successfully pass the peer review process, it is wise to consider more papers than you anticipate as the upper limit. If fewer than three articles are accepted for publication, the articles will be published as stand-alone articles in the journal.
Once you're ready, use the link below to find your chosen journal and submit your proposal.

https://www.elsevier.com/editor/role/guest/guide

\end{comment}

\begin{itemize}
\item compare VQE on classical computer vs IBM hardware to demonstrate the need
for improved hardware.
\item compare QUANDARY to other quantum optimal control algorithms. inverse
problems.
\item compare Error correction algorithms.
%4. compare emulator algorithms and compare algorithms for emulating a ``NISQC''
\item quantum and hybrid classical/quantum algorithms for benchmarking quantum 
computers.
\end{itemize}

\begin{comment}
remark:
Schrodinger equation direct model: ~$10^{3N}$  $x^t H x/x^t x$ $H$ has
dimensions of $10^{3N}$ by $10^{3N}$
Density functional theory: N * (10)
if $N$ electrons, what is the size of the Hamiltonian?
\end{comment}


\begin{comment}
The following is a template for inviting researchers to contribute to the special issue: \\
\par\noindent
Dear first and last name, \\
\par\noindent

In response to the growing interest in Quantum Computing, and the constant effort to design ever more resilient systems, we plan to publish a Special Issue on ``Numerical Methods for Engineering Quantum Computers'' in the Journal of Computational Physics (JCP). This Special Issue will span a broad range of related topics from numerical methods for determining ground states, numerical methods for solving the unsteady or nonlinear Schrodinger Equation, Density Functional Theory, Computer Aided Design of Quantum Algorithms in order to optimize the ``decoherence time,'' Hybrid classical quantum algorithms, and optimal design and control of qubits. Our special issue will serve as a complete reference on the subject matter for use by students and established researchers. \\

\par\noindent
Given your expertise in the related field, we extend a personal invitation to you to contribute a paper to this Special Issue on a topic of your choice. If agreeable, please send us a reply by Email with a tentative title by 22 November, 2024. Please copy Ms Yuan Li (yuan.li@elsevier.com) who is copied on this Email and happy to answer any questions you may have. Please see below details about the submission portal and other relevant information. \\

\par\noindent
We look forward to hearing from you. \\
\par\noindent

Our kindest regards, \\
\par\noindent

Mark Sussman, guest editor 2, guest editor 3, \ldots \\
\par\noindent

Journal:
\par\noindent

Journal of Computational Physics (ISSN: , CiteScore: , Impact Factor: ) \\
\par\noindent

Special Issue: ``Numerical Methods for Engineering Quantum Computers''
\par\noindent

Website: 
\par\noindent

Guest Editors: (e.g. 3 or more)
\par\noindent

All submissions will be peer-reviewed.
\par\noindent

Important dates:
\par\noindent

Submission Website: 
\par\noindent
select the article type of `` VSI: Numerical Methods for Engineering Quantum Computers '' 
\par\noindent

Submission portal closes: 30 November 2025.
\par\noindent

%Publication date (estimate): 30 October 2025.
Editorial acceptance deadline: 30 May 2026
\par\noindent

\end{comment}

%%\biboptions{sort&compress}
\bibliography{references}

\end{document}
%%% Local Variables:
%%% mode: latex
%%% TeX-master: t
%%% End:
